% !TeX root = ../教材题库_全部条目_转写工作稿.tex
% Source: XJ-B1-C06-TH-P251-Y0637
% 第6章 / 6.4 用样本估计总体 / 教材 p.251 / PDF 第258页
% OCR未校对：公式与文字均可能有误；未完成数学TeX转写；禁止直接发给学生
% 原题及相邻原解截图（核对用）：../题目原图/XJ-B1-C06-TH-P251-Y0637-01.webp
% 本文件只在转写完成后作为题干使用；OCR 条目尚不是可靠题干。
\begin{bankitem}{p.251 思考（含紧邻说明）}{第6章 / 6.4 用样本估计总体 / 教材 p.251 / PDF 第258页}{XJ-B1-C06-TH-P251-Y0637}{OCR未校对：公式与文字均可能有误；未完成数学TeX转写；禁止直接发给学生}
\begin{ocrdraft}
[教材p.251 / PDF第258页]
\par
？思考如何求一组观测数据的百分位数呢？
\par
首先是确定位置，然后是求出对应的百分位数．下面我们以计算尸2为例进行
\par
说明：
\par
设观测数据已经按从小到大的顺序排列，如，，\ensuremath{\cdot}
\par
第一步，计算一\ensuremath{\times}25％；
\par
第二步，如果不是整数，用m表示比c大的最小整数，则所求的P2
\par
，如
\par
》一f\ocrglyph{00AF}。十1
\par
果c是整数，则所求的P25
\par
2
\par
对于0，99〕之间的整数。将上述的25％改为厂％，即可求得P,、
\par
在统计学中，P25又称为第一四分位数，P50又称为第二四分位数，P75又称为第
\par
三四分位数．
\end{ocrdraft}
\end{bankitem}
